\chapter{Introduction}
\label{chap:introduction}

\vspace*{\fill}

Many large-scale planning problems, such as crew scheduling, vehicle routing and cutting stock,
lead to linear programs with a block structure: a number of independent subproblems, one per
crew member, vehicle or machine, are coupled by a comparatively small set of
\textbf{linking constraints}. In crew scheduling, for instance, every pilot follows a
schedule of their own, while the linking constraints require that each flight is covered by
the pilots together.

\medskip

\textbf{Column generation} is the standard method for such problems. The problem is first
rewritten in terms of complete solutions of the subproblems (a schedule, a route, a cutting
pattern), so that every variable, or \emph{column}, represents one such solution. There are
far too many columns to write down, but only a few of them are needed in an optimal
solution. Column generation therefore works with a \textbf{restricted master problem} over a
small set of columns and asks the subproblems, in the \textbf{pricing problem}, for a new
column that can improve the current solution.

\medskip

In its classical form the restricted master problem is a linear program. Its dual solutions,
which drive the pricing problem, are known to jump from one extreme point to another, which
makes the method unstable and slow in its final phase \cite{lubbecke2005selected}. A
well-established remedy is \textbf{proximal stabilization}
\cite{dumerle1999stabilized, briant2008comparison}: the dual problem is augmented by a
quadratic term that keeps the dual iterate close to a \emph{stability center}. The price of this regularization is
that every restricted master problem becomes a \textbf{convex quadratic program}.

\medskip

This work is about solving this sequence of quadratic programs efficiently, for linking
constraints that may be equalities as well as inequalities.

\vspace*{\fill}

\newpage


\section{The Computational Problem}

Two consecutive restricted master problems in column generation are almost identical. The
new problem contains one more column, and with stabilization the stability center may have
moved. A general-purpose quadratic programming solver does not exploit this: it rebuilds and
refactorizes its internal linear algebra for every master problem, and its cost per master
problem grows with the number of generated columns.

\bigskip

\textbf{Active-set methods} are well suited for sequences of closely related quadratic
programs. They move between faces of the feasible region by adding or removing one
constraint at a time, and the linear system they solve at each step changes only by a
low-rank term. In the \textbf{nullspace} variant, this system is the reduced Hessian on the
nullspace of the active constraints. Its size equals the number of degrees of freedom of the
current face, not the number of variables.

\bigskip

The question addressed in this work is how far this structure can be exploited for the
regularized master problem of column generation, in particular:
\begin{enumerate}
    \item What is the structure of the active constraints of the regularized master problem,
          and does it allow an explicit nullspace basis?
    \item How small is the linear system that has to be factorized, and how does it change
          when the working set changes?
    \item Can the factorization be carried over from one column generation iteration to the
          next, so that a new column or a new stability center costs (almost) nothing?
    \item What changes when the linking constraints are inequalities, which add sign
          conditions on the dual variables and destroy part of the structure?
    \item How does the resulting method compare in practice with general-purpose solvers and
          with classical column generation?
\end{enumerate}


\section{Contributions}

The main results of this work are the following.

\begin{itemize}
    \item \textbf{A unified formulation for mixed linking constraints.} Equality rows and
          inequality rows of both directions are written as equalities with surplus
          variables. The regularized master problem then has a single form for all cases
          (\cref{chap:formulation,chap:regularization}).

    \item \textbf{A five-set column ordering and an explicit nullspace basis.} The convexity
          constraints of the blocks, together with the surplus variables, admit an ordering
          of the variables into representatives, working-set columns, nullspace columns,
          active surpluses and free surpluses. The nullspace basis consists of vectors
          $e_j - e_{\rep(j)}$ and unit vectors, and it is never stored
          (\cref{chap:qpas,chap:nullspace}).

    \item \textbf{Decoupling and the effective reduced Hessian.} The reduced KKT system splits
          into an explicit formula for the surplus directions and a system with the
          \emph{effective reduced Hessian} $\Heff = (\Aeff Z_x)^\top(\Aeff Z_x)$. Here $\Aeff$
          holds the equality rows together with the inequality rows whose surplus is fixed at
          zero. Its dimension never exceeds the number of linking constraints
          (\cref{chap:nullspace}).

    \item \textbf{Multipliers are reduced costs.} The Lagrange multipliers of the active-set
          method are exactly the reduced costs of column generation and the dual sign
          conditions of the inequality rows. The optimality test of the master problem and
          the pricing test are therefore the same test, applied inside and outside the
          restricted master problem (\cref{chap:nullspace}).

    \item \textbf{Updates and hot start.} Every change of the working set is a row
          addition, a row deletion or a rank-one modification of a fixed-size $LDL^\top$
          factorization. A new column or a new stability center leaves the factorization
          unchanged (\cref{chap:updates}).

    \item \textbf{Exact treatment of singular reduced Hessians.} A dependent nullspace
          direction is resolved by a zero-curvature step that keeps every working-set
          variable except the released one exactly at its bound (\cref{chap:singular}).

    \item \textbf{A computational study} on generated crew scheduling instances with equality
          and inequality coverage rows: correctness against the full linear program, timing,
          the two phases of convergence, and the effect of refactorization on the fill of the
          factor (\cref{chap:results}).
\end{itemize}


\section{Outline}

\Cref{chap:formulation} introduces the multi-block linear program with mixed linking
constraints, its Dantzig--Wolfe reformulation, column generation and the surplus
formulation. \Cref{chap:regularization} adds the proximal regularization, eliminates the
auxiliary variable and derives the relation between the gradient of the quadratic master
problem and the reduced costs. \Cref{chap:qpas} sets up the active-set method for the
regularized master problem and establishes the five-set reordering of the constraint matrix.
\Cref{chap:nullspace} constructs the nullspace basis, solves the reduced KKT system with the
effective reduced Hessian and recovers the Lagrange multipliers. \Cref{chap:updates}
describes the updates of the factorization when the working set changes, the fixed-length
representation, periodic refactorization and the hot start between column generation
iterations. \Cref{chap:singular} treats singular reduced Hessians. \Cref{chap:results}
presents the computational results, and \cref{chap:conclusion} concludes.

\bigskip

Throughout, the reader is assumed to be familiar with linear programming duality, the basic
idea of column generation and the basic idea of a primal active-set method for quadratic
programming, as found in standard texts such as \cite{nocedal2006numerical, gill1981practical,
desaulniers2005column}. Everything beyond this is derived in the text.
